% Generated draft: requires mathematical and source review.
\documentclass[11pt]{article}
\usepackage[margin=25mm]{geometry}
\usepackage{fontspec}
\setmainfont{Latin Modern Roman}
\usepackage{amsmath,amssymb,mathtools,unicode-math}
\setmathfont{Latin Modern Math}
\usepackage{enumitem,needspace,xurl,hyperref}
\setlength{\parindent}{0pt}
\setlength{\parskip}{5pt}
\setlength{\emergencystretch}{3em}
\newcommand{\N}{\mathbb{N}}
\newcommand{\Z}{\mathbb{Z}}
\newcommand{\Q}{\mathbb{Q}}
\newcommand{\R}{\mathbb{R}}
\newcommand{\C}{\mathbb{C}}
\newcommand{\F}{\mathbb{F}}
\newcommand{\A}{\mathbb{A}}
\newcommand{\PP}{\mathbb{P}}
\newcommand{\E}{\mathbb{E}}
\newcommand{\eps}{\varepsilon}
\newcommand{\dd}{\,\mathrm{d}}
\begin{document}
\section*{UnsolvedMath ID: 30000746}
\textbf{Source problem code:} catalogue entry
\subsection{Definitions and mathematical statement}
For a finite simple graph $G$, the usual vertex-coloring game has Alice and Bob alternate properly coloring one uncolored vertex from a fixed palette. Alice seeks a complete coloring and Bob seeks to prevent it. With the usual Alice-first, no-pass convention, let $\chi_g(G)$ be the least palette size for which Alice can force completion.

The first question is to determine
\[
M_{\mathrm{pl}}=\sup\{\chi_g(G):G\text{ is a finite planar graph}\}.
\]
The source reports $11\leq M_{\mathrm{pl}}\leq17$ and asks for the exact extremal value. This is a bound on the maximum over the class, not a claim that every planar graph has game chromatic number at least $11$. The cited numerical range is retained as supplied historical information, not as an independently updated best bound.

The second question uses a separate online edge-coloring game. Start with an empty graph on an available supply of vertices. At each move Builder chooses a new edge, subject to the requirement that the graph of all chosen edges remains planar; Painter immediately and irrevocably colors that edge red or blue. For a fixed finite graph $H$, Builder seeks to force a monochromatic subgraph isomorphic to $H$ after finitely many moves, regardless of Painter's choices. The copy need not be induced. The proposed characterization is
\[
\text{Builder has a winning strategy}\quad\Longleftrightarrow\quad
H\text{ is outerplanar},
\]
where outerplanar means admitting a plane embedding with every vertex on the boundary of the unbounded face.

The source does not specify the first game's move-order/pass conventions, nor whether the second game's planarity constraint refers only to the current abstract graph or to a fixed embedding. The latter is normally read as abstract planarity, allowing new vertices and without prescribing a single drawing throughout. These conventions are explicit contextual interpretations, not additional verified source conditions.
\subsection{Short English statement}
Determine the largest game chromatic number among planar graphs. Separately, decide whether Builder can force a red or blue copy of a fixed graph while always keeping the board planar exactly for outerplanar targets.
\subsection{Sources}
\begin{itemize}
\item Supplied report: \emph{Mini-Workshop: Positional Games} (2008), Oberwolfach Reports, DOI 10.4171/owr/2007/20, pp. 1073--1114; proposed year 2007. \url{https://doi.org/10.4171/owr/2007/20}.
\item Local UnsolvedMath registry, permanent ID 30000746, catalogue entry; created 25 August 2026, updated 16 September 2026. The corrected\_verified statement asks for the planar maximum, resolving the original fragment's misleading per-graph wording of the $11$--$17$ range. The supplied 21 August 2026 literature triage labels both the exact-value and outerplanarity-characterization questions open. No additional external reference is supplied; neither the bounds nor current status have been independently checked here.
\end{itemize}
\end{document}
